%!TEX root = ../../main.tex
%% ===============================================================
%% 7.2 정합 입력 다층 퍼셉트론
%% 파일: chapters/07_models/02_matched_mlp.tex
%% 상위: chapters/07_models/chapter.tex
%% 정의 라벨: sec:model-mlp
%% 참조 라벨: -
%% 포함 객체: -
%% 인용: -
%% 상태: 초안 (docs/OUTLINE.md 참고)
%% ===============================================================

\section{정합 입력 다층 퍼셉트론}
\label{sec:model-mlp}

RQ3(학습기 효과)을 분리하기 위해 LightGBM과 \emph{같은} 약 2{,}980 차원 벡터를 입력으로 하는 다층 퍼셉트론(MLP)을 둔다. 이 모델은 표현을 고정한 채 학습기만 바꾼 것이므로, 두 모델의 차이는 오직 귀납 편향(트리 앙상블 대 매끄러운 함수 근사)에서 온다. 심층 하네스 안의 진단용 MLP(95 차원 거시 시퀀스의 $[x_L \,\|\, \mathrm{mean}]$ 요약 사용)는 입력이 다르므로 이 비교에 쓰지 않는다.
